%==============================================================
%  Custom commands and shortcuts
%==============================================================


%--- Frequently used physics notation -------------------------
% Note: the 'physics' package already provides \expval, \abs, \norm, \ket, ...
\newcommand{\Jrot}{J}                          % rotational quantum number
\newcommand{\Brot}{B}                           % rotational constant
\newcommand{\tHe}{\ce{^3He} }
\newcommand{\fHe}{\ce{^4He} }

\newcommand{\OCS}{\ce{OCS}\textsuperscript{\ref{tab:ocs-params} }}
\newcommand{\CSS}{\ce{CS2}\textsuperscript{\ref{tab:cs2-params} }}
\newcommand{\NOD}{\ce{(NO)2}\textsuperscript{\ref{tab:no2-params} }}

%--- Units frequently used (siunitx) --------------------------
\DeclareSIUnit{\wavenumber}{\per\centi\meter}   % cm^-1

%--- Editorial helpers (remove in final version) --------------
\newcommand{\todo}[1]{\textcolor{red}{[TODO: #1]}}
\newcommand{\ins}[1]{\textcolor{green}{[#1]}}

%--- Textbook-style aside / info box --------------------------
% Usage: \begin{aside}{Title} ... text ... \end{aside}
% Indented from the left, faint background, accent left rule,
% bold accent-coloured title line. Breaks across pages.
\newtcolorbox{aside}[1]{%
    enhanced, breakable, sharp corners,
    colback = accent!4,                              % faint background tint
    colframe = accent,
    boxrule = 0pt, leftrule = 2pt,                   % accent bar on the left only
    left = 1em, right = 1em, top = 0.7em, bottom = 0.7em,
    width = \dimexpr\linewidth-1.8em\relax,          % narrower than text...
    flush right,                                      % ...indented on the left
    before skip = 1em, after skip = 1em,
    before upper = {{\bfseries\color{accent}#1}\par\smallskip}
}
